\section{Smoothing lemmas}
\label{sec:smoothing}

\subsection{Crossings}
\label{sec:crossings}

\subsubsection{Left vs right actions of \texorpdfstring{$\pi_1(S)$}{π1(S)}}
\label{sec:actions}

To fix conventions, we briefly recall elementary facts about actions of $\pi_1(S,\ast)$ on the universal
cover~$\wt S$ of a surface (or indeed any connected length space). We
follow the convention that path composition works in the natural
order, where a product $\gamma_1 \cdot \gamma_2$ of paths in~$S$ is
defined if $\gamma_1(1) = \gamma_2(0)$.
If $a \in \pi_1(S,\ast)$, we let $[a]$ denote the free homotopy class of $a$.
In particular, we note that $[g a g^{-1}]=[a]$.

\begin{definition}[Right action]
For a covering map $p \colon \wt S \to S$ of~$S$ and base-point
$\ast \in S$, we have a \emph{right} action of $\pi_1(S,\ast)$ on
the fiber $p^{-1}(\ast)$, defined by, for $\bm{x} \in p^{-1}(\ast)$,
$\bm{x}\cdot[\gamma] = \wt \gamma(1)$, where $\wt\gamma$ is the
unique lift of~$\gamma$ with $\wt\gamma(0) = \bm{x}$.
\end{definition}
(The bold font denotes points in $\wt{S}$.)
We observe
the following.
\begin{itemize}
\item This right action does not depend on a choice of basepoint
  on~$\wt S$, and works for any covering space.%
  \footnote{One succinct statement of the classification theorem for
    covering spaces is that coverings of $S$ are in bijections with
    sets with a right action of $\pi_1(S)$.}
\item In the induced length metric on~$\wt S$, the distance between
  $\bm{x}$ and $\bm{x}\cdot [\gamma]$ is bounded by the length of~$\gamma$.
\end{itemize}

\begin{definition}[Left action]
When $\wt S$ is the universal cover of~$S$ (or, more generally, any
normal cover of~$S$), there is a \emph{left} action of $\pi_1(S,\ast)$
on all of $\wt S$ by deck transformations. If you take
the standard model of $\wt S$ where points
$\bm{P}_{[\gamma]}$ correspond
to homotopy classes of paths~$\gamma$
in~$S$ with $\gamma(0) = \ast$ (with homotopies fixing the endpoints), then this
action is defined by
$[\gamma_1] \cdot \bm{P}_{[\gamma_2]} \coloneqq
\bm{P}_{[\gamma_1\cdot\gamma_2]}$.

\DPTtodo{Change this bold-face convention}
We will also denote the point in the cover corresponding to a group
element by bold-face: $\bm{P}_{[\gamma]} = \bm{[\gamma]}$.
\end{definition}
Here we make the following observations.
\begin{itemize}
\item This left action is by an isometry of all of~$\wt S$.
\item The left action depends on a lift of
  $\ast \in S$ to a basepoint  $\bm{\wt \ast} \in \wt S$, for instance
  since the left and right actions are
  related by $[\gamma] \cdot \bm{\wt \ast} = \bm{\wt \ast} \cdot [\gamma]$.
\item The distance between $\bm{x}$ and $[\gamma] \cdot \bm{x}$
  may be unbounded, although the translation length of $[\gamma]$ on
  $\wt S$ is bounded by the length of $\gamma$ (by the first bullet
  point).
\item This left action extends to an action on the circle at infinity $\partial\wt S$.
\end{itemize}

We will frequently draw pictures of group elements $g\in\pi_1(S,\ast)$
acting on~$\bbH^2$. In these pictures, the arrows will be paths that
are lifts of~$g$, and thus correspond to the \emph{right} action on the
fiber.

For any non-trivial $x \in \pi_1(S,*)$, there is a corresponding
oriented closed curve $[x]$ on $S$, corresponding to the conjugacy
class of~$x$. We also need to specify different lifts of $[x]$ to
$\wt S$. This is conveniently specified by $x$ itself: $x$ determines
a sequence of points $\{\bm{x^k} \mid k \in \bbZ\}$, which is a
quasi-geodesic in $\wt S$, and thus determines an oriented geodesic
denoted $\ora{x}$ with endpoints $x^+ = \lim_{k \to \infty} \bm{x^k}$
and $x^- = \lim_{k \to -\infty} \bm{x^k}$. Note that $x$ acting on the
left on $\overline{\wt S}$ has north-south (NS) dynamics, with unique
attracting point $x^+$ and repelling point $x^-$: for any point
$\bm{p} \in \wt S \cup \partial \wt S$ other
than $x^-$, we have $\lim _{k \to \infty} x^k \cdot \bm{p} = x^+$.

The other lifts of $[x]$ to $\wt S$ are translates $q \cdot \ora{x}$
of $\ora{x}$ by a
group element~$q$; we can either think of this as the geodesic
$\ora{qxQ}$ passing within a bounded distance of the points $\bm{qx^kQ}$;%
\footnote{We follow the group theory convention of denoting inverses
  by a capital letter.}
or we can think of it as the geodesic corresponding to the coset
$q\langle x\rangle$, i.e., passing within a bounded distance of the
points $\bm{qx^k}$. (Note these two sequences of points are within a
bounded distance of each other, since the right action of $\pi_1$
moves points a bounded distance.)

In this way, we can see that two different lifts
\begin{align*}
  \ell_1 &= q_1 \cdot \ora{x} & \ell_2 &= q_2 \cdot \ora{x}
\end{align*}
are related by the isometry $q_2 q_1^{-1}$,
\begin{equation}\label{eq:line-translate}
  \ell_2 = q_2 q_1^{-1} \cdot \ell_1,
\end{equation}
but the geometric relation between the two lifts (up to the left
$\pi_1(S,*)$ action) is captured instead by the double coset
\begin{equation}\label{eq:double_coset}
  \langle x \rangle q_1^{-1} q_2 \langle x \rangle.
\end{equation}
Note that the order of the product is, confusingly, reversed relative to
\eqref{eq:line-translate}.

We can express many of our constructions as working with either
elements of $\pi_1(S)$ (acting by isometries on the universal cover)
or with general hyperbolic elements of the group $\Isom^+(\bbH^2)$.
(Every element of $\pi_1(S)$ necessarily acts hyperbolically on
$\wt S$.) When there is a choice, we prefer to state lemmas in terms
of $\Isom^+(\bbH^2)$, since in principle it is more general (as we do
not need to assume the group elements generate a discrete subgroup).

\DPTtodo[inline]{Deleted a bit about R-crossing here, moved it later.
  Maybe this convention is now misplaced.}
\begin{convention}\label{conv:circle-intervals}
  Intervals on the circle are always written with endpoints in counterclockwise
  order. (Thus if $x,y \in \partial \bbH^2$ are distinct points,
  $(x,y) \cup (y,x)$ is almost the entire circle).
\end{convention}

\subsection{Essential crossings}
\label{subsec:essential}
The following subsection follows closely~\cite[2.2]{MGT21:Smoothings}.

Loosely speaking, an essential crossing is a crossing of a
multi-curve that can't be homotoped away. We make this definition
precise as follows. We cover cases where $\gamma$ does not have
transverse crossings for convenience of some of the examples.

\begin{definition}[linked points on a circle]
  We say that two sets of two points $\{a,b\}$ and $\{c,d\}$ in
  $S^1$ are \emph{linked} if the four points are distinct and both
  connected components of
  $S^1- \{ a,b \}$ have an element of $\{c,d\}$.
\end{definition}

\begin{definition}[lift of a concrete curve]\label{def:lift}
  Given a concrete multi-curve~$\gamma$ on~$S$ and a choice
  $p \in X(\gamma)$, set $x=\gamma(p) \in S$. Pick a lift
  $\wt x \in \wt S$ of~$x$. The unique lifting property gives a unique
  lift $\wt{\gamma_p}: \wt{X}(\gamma;p) \to \wt S$ of $\gamma$ with
  $\wt{\gamma_p}(\wt p) = \wt x$, where $X(\gamma;p)$ is the component
  of~$X(\gamma)$ containing~$p$ and $\wt{X}(\gamma;p)$ is its
  universal cover with basepoint~$\wt p$.
\end{definition}

\begin{definition}[essential crossing]\label{def:essential-crossing}
  Let $\gamma$ be a concrete multi-curve on~$S$, and suppose we have
  points $p,q \in X(\gamma)$ so that
  $x \coloneqq \gamma(p) = \gamma(q) \in S$. Pick a lift
  $\wt x \in \wt{S}$ of $x$, and let $\wt{\gamma_p}$ and
  $\wt{\gamma_q}$ be
  the corresponding lifts of components of~$\gamma$ following
  Definition~\ref{def:lift}.
  Then the pair $(p,q)$ form an \emph{essential crossing} if the following two conditions hold:
  \begin{enumerate}
  \item both components of $X(\gamma)$ containing $p$ and $q$ are not
    null-homotopic, so that $\wt{\gamma_p}$ and $\wt{\gamma_q}$ are quasi-geodesic components of $\wt \gamma$; and
  \item either
    \begin{enumerate}
    \item the endpoints $\{a,b\}$ of $\wt{\gamma_p}$
      and the endpoints $\{c,d\}$ of $\wt{\gamma_q}$
      are linked in $S^1_\infty$; or
    \item\label{item:ess-multi} $p$ and $q$ lie on the same component of $X(\gamma)$,
      $[\gamma] = [\delta^n]$ for some $n > 1$ and some
      primitive $\delta \in \pi_1(S,x)$,
      the loop from $p$ to~$q$ in $X(\gamma)$ maps to $[\delta^m]$ for
      some $0 < m < n$, and so the loop from $q$ to~$p$ maps to $[\delta^{n-m}]$.
    \end{enumerate}
  \end{enumerate}
\end{definition}
In case~\ref{item:ess-multi}, the endpoints of $\wt {\gamma_p}$ and
$\wt{\gamma_q}$ are the same.

This definition might be somewhat looser than expected. For
instance, in the chain of three crossings
\[
  \mfig{curves-20}
\]
the circled middle crossing is
essential iff the other two are, since ``linking at infinity'' doesn't
see the direction of crossing. This does not matter for our purposes.

\begin{definition}[smoothings]
  \label{def:smoothing}
  Let $(p,q) \in X(\gamma)$ be an essential crossing of~$\gamma$
  on~$S$. To make a \emph{smoothing}~$\gamma'$ of $(p,q)$, cut
  $X(\gamma)$ at $p$ and~$q$ and reglue the resulting four endpoints
  in one of the two other possible ways, getting a new $1$-manifold
  $X(\gamma')$. The map $\gamma'$ agrees with~$\gamma$; this is
  well-defined since $\gamma(p) = \gamma(q)$. In pictures we will
  homotop $\gamma'$ slightly to round out the resulting
  corners. If $\gamma$ is oriented, then the \emph{oriented smoothing}
  is the smoothing that respects the orientation on $X(\gamma)$:
  \begin{equation*}
    \mfig{curves-3} \reducesto \mfig{curves-4}
  \end{equation*}
If we obtain a concrete curve $\gamma'$ from $\gamma$ by a sequence of
$k$ smoothings of essential crossings,  we will write
$\gamma \reducesto_k \gamma'$. (We check whether the
crossings are essential at each stage of this process; this is
more restrictive than checking at the beginning.)
\end{definition}

\begin{lemma}[{\cite[Lemma~2.8]{MGT21:Smoothings}}]
  Essential crossings are unavoidable in a homotopy class, in the
  sense that if $\gamma$ and $\gamma'$ are homotopic concrete
  multi-curves and $(p,q) \in X(\gamma)$ is an essential crossing of~$\gamma$,
  then there is an essential crossing $(p',q') \in X(\gamma')$ so
  that the smoothings of $(p,q)$ and of $(p',q')$ are homotopic.
\end{lemma}

We will distinguish two types of smoothings.

A \emph{disconnected smoothing} is a smoothing of a crossing of two distinct curves.
A \emph{connected smoothing} is a smoothing of a self-crossing of two distinct curves.

\subsubsection{Right-handed crossings}
\label{subsec:rhcrossings}
In order to work better with group elements and crossings between
oriented geodesics, we pick out one type of essential crossing.

\begin{definition}[Right handed crossing]
  \label{def:righthandedcrossing}
  For non-trivial $x,a \in \pi_1(S)$ that do not share a common power,
  the corresponding geodesics $\ora{a}$ and $\ora{x}$ have a
  \emph{right-handed crossing} or \emph{R-cross} if the endpoints
  $a^-, x^-, a^+, x^+$ appear in that
  (counter-clockwise) cyclic order on $\partial \wt S$. (We also apply
  rules of English grammar, for instance talking about left-handed
  crossings if they have the opposite cyclic order.)
\end{definition}

Note that, if we draw the loops $[x]$ and $[a]$ as based at the
basepoint $* \in S$, the loops intersect at $*$ by definition. If that
intersection is essential (and $x$ and $a$ do not share a power),
there are two possible cyclic orders for the endpoints, so either
$\ora{a}$ crosses $\ora{x}$ to the right or $\ora{x}$ crosses
$\ora{a}$ to the right. In the first case (as in the Definition) the picture looks
intuitively like
\[\mathcenter{\resizebox{20mm}{!}{\svgfig{rhres}}}.\]

\DPTtodo{This comment deserves a picture}
If we wish to talk about some essential crossing of $[a]$ and $[b]$
that is \emph{not} at the basepoint, we can always conjugate by a
suitable element~$c$, and ask about the crossing between $\ora{caC}$
and $\ora{b}$. Concretely, to find $c$, homotop $[a]$ and $[b]$ so
they pass through the basepoint~$*$, and take $c$ to be the homotopy
class of the loop that runs along~$[b]$ from $*$ to the chosen
intersection, and back along~$[a]$ from the chosen intersection
to~$*$.

Since a right-handed crossing is essential, we have the following.
\begin{lemma}
Let $a,x \in \pi_1(\Sigma)$. If $\ora{a}$ crosses $\ora{x}$ to the
right (or to the left), then
\[
  [a][x] \reducesto [ax]\qquad\text{and}\qquad[a][x] \reducesto [aX].
\]
\end{lemma}

Observe that whether a crossing is ``right-handed'' is invariant under
simultaneous conjugation of $a$ and $x$, but does depend on
more than the conjugacy classes of $a$ and~$x$. (We need to specify
which crossing we are considering.)

In terms of smoothings, in a disconnected smoothing, $a$ and $x$ above are non-conjugates and non-inverses of each other. In a connected smoothing, 
$a$ and $x$ above are conjugate.


We also make the following trivial observation, using
Convention~\ref{conv:circle-intervals}.
\begin{lemma}\label{lem:hyperbolic}
  The geodesics $(\ora{a}, \ora{x})$ R-cross iff
  $a \cdot (x^-, x^+) \subsetneq (x^-, x^+)$. More generally, if
  $a \in \Isom(\bbH^2)$ is a hyperbolic element with attracting
  endpoint $a^+$ in an interval $(x,y)$ iff
  $a \cdot (x,y) \subsetneq (x,y)$.
\end{lemma}

% Further, it is anti-symmetric: if
% the crossing between $a$ and~$x$ is right-handed, then the crossing
% between $x$ and~$a$ is left-handed.



\subsubsection{Parallelism and crossing of hyperbolic axes}

We now turn to some elementary lemmas about crossing. For this
purpose, we also distinguish the various other types of relations
between two geodesics $\ora{a}, \ora{b}$.

We will switch to using elements of $\Isom^+(\bbH^2)$. On the one
hand, if we fix any hyperbolic metric on~$S$, any non-trivial element
$a \in \pi_1(X)$ acts on $\wt S = \bbH^2$ on the left by an isometry
of $\bbH^2$ that is hyperbolic (in the sense of having two distinct
fixed points on $\partial\bbH^2$). On the other hand, the notions we
will consider will only depend on the cyclic order on the circle at
infinity of endpoints of geodesics, and in particular do not depend on
the metric.
In principle one should invoke a representation $\rho \colon \pi_1(S)
\to \Isom^+(\bbH^2)$, but since nothing depends on this representation
we will suppress it.
% \DMGtodo[inline]{(Punctures) If punctures, we need to ask for
%   hyperbolic, and peripheral sent to parabolic.}

\begin{definition}
  Let $a,b$
  be hyperbolic elements in $\Isom^+(\mathbb{H}^2)$, and let
  $\ora{a}, \ora{b}$ be the oriented geodesics that are the
  corresponding axes.
\begin{itemize}
    \item As before, the pair $(\ora{a},\ora{b})$ \emph{cross to the
        right} or \emph{R-cross} if $a^-, b^-, a^+,b^+$ appear in this
      ccw order on the boundary at infinity. See
      Figure~\ref{fig:cross}.
    \item The pair $(\ora{a},\ora{b})$ is \emph{R-parallel} if $b^-,b^+, a^+, a^-$ appear in this ccw order on the boundary at infinity. See Figure~\ref{fig:parallel}.
   \item  The pair $(\ora{a},\ora{b})$ is \emph{R-anti-parallel} if $b^+,b^-, a^+, a^-$ appear in this counter-clockwise order on the boundary at infinity. See Figure~\ref{fig:anti}.
   \end{itemize}
   We can similarly define L-cross, L-parallel, and L-anti-parallel.
\end{definition}

If we swap $\ora{a}$ and $\ora{b}$, the type of crossing/parallelism
can change, as explained below.

\begin{lemma} Let $a,b\in \Isom^+(\bbH^2)$ be hyperbolic elements.
\begin{enumerate}
  \item If $(\ora{a},\ora{b})$ R-cross, then $(\ora{A},\ora{B})$
    R-cross and $(\ora{b},\ora{a})$ L-cross.
  \item If $(\ora{a},\ora{b})$ are R-parallel, then
    $(\ora{B},\ora{A})$ are R-parallel, $(\ora{b},\ora{a})$ are
    L-parallel, and $(\ora{a},\ora{B})$ are
    R-anti-parallel.
  \item If $(\ora{a},\ora{b})$ are R-anti-parallel, then
    $(\ora{b},\ora{a})$ are R-anti-parallel, $(\ora{A},\ora{B})$ are
    L-anti-parallel,
    and $(\ora{a},\ora{B})$
    are R-parallel.
\end{enumerate}
\end{lemma}

\begin{proof}
  Immediate from the definitions (and this lemma can be taken to be a
  definition of the ``L'' variants).
\end{proof}

\begin{lemma}
  Let $a,x \in \Isom^+(\bbH^2)$ be hyperbolic elements that do not share any endpoints.
  Then $\ora{a}$ and~$\ora{x}$ R-cross (resp.\ L-cross) iff
  $(\ora{x},a \cdot \ora{x})$ are R-parallel (resp.\ L-parallel).
 \label{lem:rhs}
\end{lemma}
\begin{proof}
  The statement that $(\ora{x},a\cdot\ora{x})$ are R-parallel is (by
  definition) equivalent to $a\cdot(x^-,x^+) \subsetneq (x^-,x^+)$,
  which in turn is equivalent to $a^+ \in (x^-,x^+)$ and
  $a^- \notin (x^-,x^+)$, which is turn is equivalent to
  $(\ora{a},\ora{x})$ R-crossing.
\end{proof}

We also have the
following. 

\begin{lemma}
  Let $a,b,x \in \pi_1(S)$ be group elements so that $(\ora{a},\ora{x})$ and $(\ora{b},\ora{x})$ R-cross. Then:
  \begin{enumerate}
  \item All four geodesics $\ora{xa}$, $\ora{Xa}$, $\ora{ax}$, and
    $\ora{aX}$ R-cross $\ora{x}$.
  \item $(\ora{ab},\ora{x})$ R-cross.
  \item More generally, for any $n,k \in \bbZ$, $(\ora{a x^n},\ora{x}$
    and $(\ora{ax^nbx^k},\ora{x})$ R-cross.
  \end{enumerate}
    \label{lem:rhs_prod}
\end{lemma}
\begin{proof}
  For the first claim, consider the action of the respective group
  element on the interval $(x^-,x^+)$: we have
  $x^{\pm1} \cdot (x^-,x^+) = (x^-,x^+)$ and
  $a \cdot (x^-,x^+) \subsetneq (x^-,x^+)$, so four of $ax^{\pm1}$ and
  $x^{\pm1}a$ map this interval strictly inside itself. It follows
  that all these elements are hyperbolic and R-cross $\ora{x}$.

  For the second claim, we similarly have that both $a$ and $b$ map
  $(x^-,x^+)$ strictly inside itself, so their product does as well.

  The last claim is an immediate consequence of the first two parts.
\end{proof}


\begin{lemma} \DMGtodo[inline]{This lemma might be combined in the future with the more precise statements about smoothings in the application section}
  Let $a,b\in \Isom^+(\bbH^2)$ be hyperbolic elements. Then:
  \DPTtodo[inline]{Maybe we need statements about placement of
    attracting/repelling endpoints. If so, add into the cases}
  \begin{enumerate}[label=(\Roman*)]
  \item\label{item:ab-cross} If $(\ora{a}, \ora{b})$ R-cross, then $ab$ and $aB$ are also
    hyperbolic, and the pairs
    $(\ora{ba},\ora{ab})$ and $(\ora{aB},\ora{Ba})$ are both R-parallel.
  \item\label{item:ab-parallel} If $(\ora{a},\ora{b})$ are R-parallel, then $ab$ is hyperbolic
    and $(\ora{ba},\ora{ab})$ R-cross.
  \item If $(\ora{a},\ora{b})$ are R-anti-parallel, then $aB$ is
    hyperbolic and $(\ora{Ba},\ora{aB})$ R-cross.
  \end{enumerate}
  \label{lem:righthandedfoundation}
\end{lemma}


\begin{proof}
  Note that the actions of $a$ and $b$ on geodesics act by conjugation
  on the underlying word. Specifically, this means that we have
  \[
    \begin{tikzpicture}[x=1.5cm]
      \node (ba) at (0,0) {$\ora{ba}$};
      \node (ab) at (1,0) {$\ora{ab}$};
      \node (Ba) at (3,0) {$\ora{Ba}$};
      \node (aB) at (4,0) {$\ora{aB}$.};
      \draw [->,bend left=15] (ba) to node[above,cdlabel]{a} (ab);
      \draw [->,bend left=15] (ab) to node[below,cdlabel]{b} (ba);
      \draw [->,bend left=15] (Ba) to node[above,cdlabel]{a} (aB);
      \draw [->,bend right=15] (Ba) to node[below,cdlabel]{b} (aB);
    \end{tikzpicture}
  \]
  The hyperbolic elements $a$ and~$b$ have north-south dynamics. Using
  this in the three cases depending on the relation of $\ora{a}$ and
  $\ora{b}$ (crossing, parallel, or anti-parallel), we have the
  following.
  \begin{enumerate}[label=(\Roman*)]
  \item If $(\ora{a},\ora{b})$ R-cross, then we claim that the
    endpoints of the various geodesics must be laid out as in
    Figure~\ref{fig:cross}.
    \DPTtodo{Combine Figs. 3.1--3.3 into one figure/subfigures. Mark
      the $a$ and $b$ arrows to
      indicate group actions. Captions are outdated.}
    \begin{figure}
\centering{
\resizebox{60mm}{!}{\svgfig{rcrossing}}
    \caption{$\ora{a}$ and $\ora{b}$ R-cross, together with the other crossing relations for related words. }
    \label{fig:cross}}
    \end{figure} 
    To see this, observe that, by the dynamics of $a$ and $b$ on the
    circle as shown, $ab$ maps the interval
    $(a^+,b^+)$ strictly inside itself, which implies that it is
    hyperbolic with attracting endpoint $(ab)^+$ in that interval;
    similarly we see $(ab)^- \in (a^-,b^-)$. The actions of
    $a$ and $b$ interchange $(ab)^- \leftrightarrow (ba)^-$ and
    $(ab)^+ \leftrightarrow (ba)^+$, forcing the configuration shown.
    (Note that all these endpoints $(ab)^\pm,(ba)^\pm$ are on the
    portions of the circle
    where $a$ and $b$ move in opposite directions.)

    Since in this situation $(\ora{B},\ora{a})$ also R-cross, the
    other conclusion follows.
  \item If $(\ora{a},\ora{b})$ are R-parallel, then again the dynamics
    imply that $ab \cdot (b^+,a^+) \subsetneq (b^+,a^+)$, so $ab$ is
    hyperbolic and $(ab)^+ \in (b^+,a^+)$. The same interchange
    property then forces
    $(ab)^+, (ba)^+ \in (b^+,a^+)$  and
    $(ab)^-, (ba)^- \in (a^-,b^-)$. From this we see that
    $(\ora{ba},\ora{ab})$ R-cross, as in Figure~\ref{fig:parallel}.
    \DPTtodo{Add direction arrows for action of $a,b$ to the illustration}

    \begin{figure}
\centering{
\resizebox{60mm}{!}{\svgfig{parallel}}
    \caption{$a$ and $b$ are R-parallel, together with the other crossing relations for related words. }
    \label{fig:parallel}}
    \end{figure}
  \item If $\ora{a}$ and $\ora{b}$ are anti-parallel, we can reduce to
    the parallel case by replacing $b$ with $B$, arranged as in
    Figure~\ref{fig:anti}.\qedhere
  \end{enumerate}
\end{proof}
    \begin{figure}
\centering{
\resizebox{60mm}{!}{\svgfig{antiparallel}}
    \caption{$a$ and $b$ are anti-parallel, together with the other crossing relations for related words. }
    \label{fig:anti}}
    \end{figure}

    See Section~\ref{sec:appsintersection} for a more in depth analysis of the configurations of lifts that can arise from disconnected smoothings.

\begin{remark}
  In the R-parallel case we do not make claims about the
  relationship of $\ora{aB}$ and $\ora{Ba}$. Indeed this is not
  determined by the hypotheses.  Similarly the R-anti-parallel case says nothing about the relation
  of $\ora{ab}$ and $\ora{ba}$.
  See Section~~\ref{subsec:connected_smoothing}, in particular Figure \ref{fig:r3} and Figure \ref{fig:r1} for the full analysis.
\end{remark}

  Lemma~\ref{lem:righthandedfoundation} applies to disconnected
  smoothings. For a similar result for a
  (self-)crossings of $a$ $\ora{ab}$ with $\ora{ba}$ one has to deal with several cases  (This is related to the
  previous remark.) In particular if $\ora{ab}$ R-crosses $\ora{ba}$,
  then $\ora{a}$ and $\ora{b}$ could be parallel or anti-parallel. The
  full analysis is similar to the discussion of maximal bracelet
  chains in the second author's earlier work \cite{Thurston14:PBI}.

\subsection{Connected smoothings}

If $C$ represents a curve with a self-intersection decomposing as $ab$, we can consider intersecting axes $\ora{ab}$ and $\ora{ba}$ in the universal cover so that their intersection is a lift of the self-intersection of the curve $ab$.

One of the resolutions of $[ab]$ is the multi-curve $[a] \cup [b]$.
The lifts of $[a]$ and $[b]$ to the universal cover correspond to axes $\ora{a}$ and $\ora{b}$ which are parallel. \DMGtodo[inline]{does this need justification? This is essentially saying that we only need to consider the case $a$ and $b$ parallel, and we do not need to worry about $a$ and $b$ being anti-parallel.}
It remains to determine the position of the lifts $\ora{aB}$ and $\ora{Ba}$.

To that end, consider the common perpendicular line of $a$ and $b$. Let $R_1$ be the reflection along this line.
Let $R_2, R_3, R_4, R_5$ be reflections so that $a=R_1 R_2 = R_4 R_1$, and $b=R_1 R_3 = R_5 R_1$. Note that $R_i$ are reflections along pairwise non-intersecting reflection lines, since $aB$ and $ab$ are hyperbolic.
It then follows that 
$ab=R_4 R_3$ and $ba=R_5 R_2$. Furthermore, $Ba=R_3 R_2$ and $aB=R_4 R_5$.
We have the following possible four configurations of non-intersecting reflection lines. 

\begin{figure}
     \centering
     \begin{subfigure}[b]{0.45\textwidth}
\resizebox{80mm}{!}{\svgfig{r3}}
    \caption{Case 1}\label{fig:r3}
     \end{subfigure}
     \hfill
     \begin{subfigure}[b]{0.45\textwidth}
\resizebox{80mm}{!}{\svgfig{r2}}
    \caption{Case 2 (not possible)}\label{fig:r2}
\end{subfigure}\\
\begin{subfigure}{0.45\textwidth}
     \resizebox{80mm}{!}{\svgfig{r1}}
     \caption{Case 3}\label{fig:r1}
   \end{subfigure}\hfill
\begin{subfigure}{0.45\textwidth}
  \centering{\resizebox{45mm}{!}{\svgfig{c2}}}
    \caption{In case 3, $aB$ and $Ba$ can cross.}\label{fig:c2}
   \end{subfigure}\hfill
   \caption{Topological cases for
     connected smoothing of $\ora{ab}$ with $\ora{ba}$.}
\end{figure}
\DPTtodo{Combined figures, but they are not uniformly cropped to same
  size. In general it's better to scale by a percentage than to an
  absolute width}

In the figures, we also draw their corresponding translation axes for each isometry.
We observe that case 2 can be ruled out, since the NS dynamics of $b$ should translate $\ora{Ba}$ to $\ora{aB}$, and thus these two axes should stay on the same side of $\ora{b}$, which does not happen in that configuration.
We note there is another case, analogous to case 2, where the reflection lines below $R_1$ are exchanged by the ones above it, which can be similarly ruled out.

We observe that in case 3, $\ora{aB}$ and $\ora{Ba}$ may cross (Figure~\ref{fig:c2}).



\subsection{Fundamental smoothing lemmas}

As part of our construction, we will use large powers of some fixed
elements. We collect some propositions guaranteeing essential
crossings in this context.
% \begin{figure}
% \centering{
% \resizebox{90mm}{!}{\svgfig{mn_plane}}
% \caption{The smoothings covered by
%   Proposition~\ref{prop:join-split-geom}, shown in the $(m-n,r)$
%   plane. Allowed regions are shaded in grey.}
% \label{fig:mn_plane}
% }
% \end{figure}

\begin{proposition}\label{prop:join-split-geom}
Let $a,b,x \in \pi_1(S)$ be arbitrary curves, with $a$ and $b$
crossing $x$ to the right. Then there is $s \ge 0$ (depending on $a,b$)
  so that, for all
 $n, m, r \in \mathbb{Z}$ with either $s < r < m-n-s$ or $m-n+s < r <
 -s$, we have
 \DPTtodo{May have messed up the statement in the editing}
  \begin{align}
    \label{eq:cross_a} [a x^n] [b x^m] &\reducesto [ax^{n+r}bx^{m-r}]\\
    \label{eq:cross_b} [a x^nb x^m] &\reducesto [ax^{n+r}][bx^{m-r}].\\
    \intertext{There exists $s' > 0$ (depending on $a$) so that for
    all $n,m$ with $|n-m| \ge s'$,}
    [ax^n][ax^m] &\reducesto [x^{\abs{n-m}}].\label{eq:cross-red3}\\
    \intertext{If $x$ is simple and $a=b$, we may take $s=0$ and
    $s'=1$. In particular, for all $n\in\bbZ$,}
    [a x^{n+1}] [a x^{n-1}] &\reducesto 2[ax^{n}].\label{eq:cross-red2}
  \end{align}
\end{proposition}
% We will also use this lemma in Section~\ref{sec:recovery} to show
% $\mu_f$ recovers the value of $f$ on curves, so we consider cases when
% $x$ is not simple.

For the allowed values of $r$, Eqs.\ \eqref{eq:cross_a} and
\eqref{eq:cross_b} move the exponents $n,m$ closer to the diagonal
$m=n$.

\begin{proof}
  In each case, we must show that we can choose lifts for which a
  relevant crossing is essential. We use the bands model, centered
  on~$\ora{x}$, and normalized so that $x$ acts by translation to the
  right by $1$. We will
  consider intervals on $\bdy^\infty\wt S$ cut out by certain other
  lifts of~$x$. Specifically, let $w(x^ka) = (x^ka\cdot x^-, x^ka\cdot
  x^+)$ be the interval on $\bdy^\infty \wt{S}$ cut out by the
  endpoints of $x^k a \cdot \ora x$
  that does not contain the endpoints $x^\pm$ of $\ora x$. Similarly
  for $w(x^k a^{-1}) = x^ka^{-1} \cdot (x^+,x^-)$,
  $w(x^k b) = x^k b \cdot (x^-,x^+)$, and
  $w(x^k b^{-1}) = x^kb^{-1}\cdot (x^+,x^-)$. All of
  these are intervals
  on the upper or lower boundary in the bands model, by the assumption
  that $a$ and~$b$ cross $x$ to the right.

  \DPTtodo{Expanded this key fact}

  Furthermore, $\ora{x^kax^\ell}$ R-crosses $x^ka \cdot \ora x$. To
  see this, conjugate by $(x^k a)^{-1}$: this is equivalent to
  $\ora{a x^{k+\ell}}$ R-crossing $\ora x$, which is immediate from
  Lemma~\ref{lem:rhs_prod}. Thus $(x^ka x^\ell)^+ \in w(x^k a)$.
  Similarly, $\ora{x^k a x^\ell}$ R-crosses $x^{-\ell}a^{-1} \cdot
  \ora{x}$ (which we can see by conjugating by $a x^\ell$), so
  $(x^k a x^\ell)^- \in w(x^{-\ell a})$.
  See Figure~\ref{fig:abcrossings}.

\DPTtodo{Figure \ref{fig:abcrossings} has $a$ crossing $x$ to the
  left. Reflect it vertically. Also some labels are misplaced}
  \begin{figure}
\centering{
\resizebox{150mm}{!}{\svgfignormal{abliftsbands}}
\caption{The bands model centered on $\ora{x}$ used in
  Proposition~\ref{prop:join-split-geom}, in a case where $[x]$ is simple.
  Lifts of $x$ are depicted in blue. The central blue lift relates to
  the lower and upper lifts by translates of the type $x^k a$, $x^l
  a^{-1}$, $x^m a$ and $x^n a^{-1}$, as labelled. Each translate
  determines a window $w( \cdot )$, which is an interval on
  $\partial_{\infty} \Sigma$ determined by its endpoints at infinity.
  If $[x]$ is not simple, the blue lifts can intersect each other.}
\label{fig:abcrossings}
}
\end{figure}


  Observe that (for instance) $w(x^n a)$ is a translate of $w(a)$.
  Furthermore, each window is a
  finite size, so there is a finite range $[s_1,s_2]$ so that $w(a)$
  is to the left of (resp.\ to the right of) $w(x^rb)$ if $r < s_1$
  (resp.\ $r > s_2$). Similarly there
  is a range $[s_1',s_2']$ so that $w(a^{-1})$ is disjoint from
  $w(x^rb^{-1})$ if $r \notin [s_1',s_2']$. Set
  $s \coloneqq \max\{s_2,s_2',-s_1,-s_1'\}$.

  Now consider the curves 
  $\ora{ax^n}$ and $\ora{x^rbx^{m-r}}$, with $s < r < m-n-s$. We have
  \[
    r-m < -n-s \le -n+s_1',
  \]
  so $w(x^{r-m}b^{-1})$ is to the left of $w(x^{-n}a^{-1})$, and
  \[
    s_2 \le s < r,
  \]
  so $w(x^rb)$ is to the right of $w(a)$. Combined, these mean that
  $x^rbx^{m-r}$ crosses $ax^n$ to the left. Similarly if
  $m-n+s < r < -s$, the curve $x^rbx^{m-r}$ crosses $ax^n$ to the
  right. In either case, we have
  \[
    [ax^n][bx^m] = [ax^n][x^r bx^{m-r}] \reducesto [ax^{n+r}bx^{m-r}]
  \]
  proving Equation~\eqref{eq:cross_a}.

  For Equation~\eqref{eq:cross_b}, write $ax^nbx^m$ as $gh$ with
  $g=ax^{n+r}$ and $h=x^{-r}bx^m$. Note the endpoints of $gh$
  are in $w(x^{-m}b^{-1})$ and $w(a)$, while the endpoints of its
  conjugate $hg =x^{-r}bx^max^{n+r}$ are in $w(x^{-n-r}a^{-1})$ and
  $w(x^{-r}b)$. Then if $s < r < m-n-s$, we have
  \[
    -r < -s \le -s_1,
  \]
  so $w(x^{-r}b)$ is to the left of $w(a)$, while
  \[
    -m \le -m+s+s_1' < -n - r + s_1',
  \]
  so $w(x^{-n-r}a^{-1})$ is to the right of  $w(x^{-m}b^{-1})$, and
  $gh$ has an essential self-crossing giving Eq.~\eqref{eq:cross_b}.
  We similarly get an essential self-crossing if $m-n+s < r < -s$.

  In the cases when $x$ is simple, two windows are disjoint, nested,
  or identical; this implies that $s_1 = s_2$ and $s_1' = s_2'$. If in
  addition $a=b$, we also see that all the $s_i$ are equal to~$0$.
  Eq.~\eqref{eq:cross-red2} follows immediately.

  To see Eq.~\eqref{eq:cross-red3}, suppose without loss of generality
  that $n < m$. Let $s_0(a) \coloneqq \lfloor \wid{x}(w(a)) \rfloor$.
  Observe that, since the bands model is normalized so that the
  translation length of $x$ is~$1$, if $x$ is simple, necessarily
  $\wid{x}(w(a))<1$ (for otherwise two lifts of $x$ would intersect),
  and thus in that case $s_0=0$. Assume $n-m > s_0$.  Then by the
  analysis above, $(ax^n)^-$ is in $w(x^{-n}a^{-1})$, and thus to the
  right of $(ax^m)^-$, which is in $w(x^{-m}a^{-1})$. For the forward
  endpoints, the above simple analysis just says that both endpoints
  lie in $w(a)$. We can do better by looking deeper in the words:
  $(ax^n)^+$ lies in $w(ax^na)$, meaning the interval between the
  endpoints of $ax^na\cdot\wt x$ that is disjoint from $x^\pm$.
  Similarly $(ax^m)^+$ lies in $w(ax^ma)$, which is to the right of
  $w(ax^na)$. It follows that $ax^n$ crosses $ax^m$ to the right, and
  so
  \[
    [ax^n][ax^m] \reducesto [(ax^n)^{-1}ax^m] = [x^{m-n}]
  \]
  as desired.
\end{proof}

\begin{remark}
  In the proof of Proposition~\ref{prop:join-split-geom}, the only
  place we used an unoriented smoothing was in the proof of
  Eq.~\eqref{eq:cross-red3}.
  In fact, this is the only place we use unoriented smoothings in the paper.
\end{remark}

We will also need the following converse smoothing.

\DPTtodo{Put this in the bigger lemma}
\begin{lemma}
Let $a,x \in \pi_1(\Sigma)$ so that $a$ crosses $x$ to the right. There exists $s_0=s_0(a) \in \mathbb{N}$ so that
\[
[ax^n][a] \reducesto [x^n], \mbox{ for all } n \geq s_0,
\]
\label{lem:Q}
\end{lemma}
\begin{proof}
The claim follows by Equation~\ref{eq:cross-red3} in Proposition~\ref{prop:join-split-geom}, taking $m=0$.
\end{proof}

We also need some lemmas about smoothings where we have high powers of
two different elements.

\DPTtodo{I don't think this lemma belongs in this section at all, move
  it to where we use the lemma}
\begin{lemma}
  Given $x, y \in \pi_1(S)$, there exist $a,b \in \pi_1(S)$ so that
  $b^{-1} \cdot \ora{x}$, $\ora{y}$, $a \cdot \ora{x}$ and
  $ab \cdot \ora{y}$ are R-parallel, in this order, as in
  Figure~\ref{fig:initial}.
\label{lem:initial}
\end{lemma}
\begin{proof}
  Given $y$, by Lemma~\ref{lem:diagonaldense}, we can find a translate
  of $\ora{x}$, denoted $b^{-1} \cdot \ora{x}$, so that
  $b^{-1} \cdot \ora{x}$ and $\ora{y}$ and are R-parallel. For any $a$,
  be translation we see $a \ora{x}$ and $ab \cdot \ora{y}$ are
  R-parallel. Find $g \in \pi_1(X)$ so that $g^{+}\in(y_-, y_+)$ and
  $g^-\in((b^{-1} \cdot \ora{x})^+, (b^{-1} \cdot \ora{x})^-)$. Set
  $a = g^n$ for large~$n$. By the North-South
  dynamics of $g$, for $n$ is large enough $a \cdot \ora{x}$ is R-parallel to $\ora{y}$, as
  desired.
\end{proof}

\begin{figure}
\centering{
\resizebox{60mm}{!}{\svgfig{initial}}
\caption{Setup configuration for fundamental smoothing lemma.}
\label{fig:initial}
}
\end{figure}

\begin{proposition}\label{prop:fundamental-aybx}
Let $a,b,x,y \in \pi_1(S)$ be arbitrary curves, ordered as in Figure~\ref{fig:initial}.
There exist $r_0,s_0$ so that for all $r > r_0$ and $s>s_0$, we have

  \begin{align}
    \label{eq:crossuniform_n} \mbox{For all } m \in \mathbb{Z}, \mbox{ } [b y ^m a x^{n+r}] [b y^m a x^{n-r}] &\reducesto [(b y^m a x^{n})^2]\\
    \label{eq:crossuniform_m} \mbox{For all } n \in \mathbb{Z},\mbox{ } [b y^{m+s} a x^{n} ] [b y^{m-s} a x^{n}] &\reducesto [(b y^m a x^{n})^2]
  \end{align}
\end{proposition}


\DPTtodo{Figure~\ref{fig:aybx} is too busy, I think there are too many
  lines shown and it's not clear what the key parts are.}
\begin{figure}
\centering{
\resizebox{170mm}{!}{\svgfignormal{aybx}}
\caption{Bands model centered on $\ora{y}$. Translates of $\ora{x}$ are coded in blue, while translates of $\ora{y}$ are coded in red. The translates of $\ora{x}$ by $y^{-s} b^{-1}$ and $b^{-1}$ are shown and labelled in the upper part, and translates of $\ora{x}$ by $y^{m-s} a$ and $y^m a$ are shown in the lower part. }
\label{fig:aybx}
}
\end{figure}

\DPTtodo[inline]{Following lemma misplaced, have not yet introduced window}
The following lemma follows from Lemma~\ref{lem:rhs}.

\begin{lemma}
Given $a,b, x \in \pi_1(\Sigma)$, so that $b^{-1} \cdot \ora{x}$ and $a \cdot \ora{x}$ are parallel, for every $n \in \mathbb{Z}$, $(ax^n b)^+ \in w(a)$ and $(ax^n b)^- \in w(b^{-1})$.
\label{lem:cornering}
\end{lemma}




\begin{proof}[Proof of Proposition~\ref{prop:fundamental-aybx}]
  We prove Equation~\ref{eq:crossuniform_m} by showing  choose lifts of conjugates of the curves in the left hand-side of the smoothing for which the crossing is  essential. Use the band model, centered on~$\ora{y}$, and normalized so that the translation length of $y$ along $\ora{y}$ is equal to 1. This is what we will mean by \emph{normalized bands model}.  \DMGtodo[inline]{We use this several times, we should extract a definition} We will
  consider the windows on $\bdy^\infty\wt \Sigma$ cut out by certain other
  lifts of~$x$. Specifically, let $w(y^m a)$ be the interval on the
  boundary at infinity cut out by the endpoints of $y^m a \cdot \wt x$
  that does not contain the endpoints $y^\pm$ of $\wt y$. Similarly for
  $w(y^{m-s} a)$, $w(y^{-s} b^{-1})$, and $w(b^{-1})$. All of these are intervals
  on the upper or lower boundary in the bands model, by the initial configuration assumption of Lemma~\ref{lem:initial}.
  See Figure~\ref{fig:aybx}; ``left'' and ``right'' will refer
  to the order on the boundary (minus $x^\pm$) of that figure.

  Observe that (for instance) $w(y^{m} a)$ is a translate of $w(a)$.
  Furthermore, each window is a
  finite size, so there is a number $s>0$ so that $w(y^{m-s}a)$
  is to the left of $w(y^m a)$, and similarly,
   $w(y^{-s}b^{-1})$ is to the left of
  $w(b^{-1})$.
  In fact, since the translation length of $y$ is 1, it suffices to take
  \begin{equation*}
  s \geq s_0 = \lfloor \max \{ \wid{y}(w(a^{-1})),\wid{y}(w(b)) \} \rfloor
  \end{equation*}
  where $\wid{y}(w(b))$ denotes the length of the interval defining the window $w(a)$ at the boundary at infinity of the normalized bands model, where distances in the boundary are measured with the Euclidean distance.
  Consider now the curves
  $y^{m-s}ax^nb$ and $y^m ax^n by^s \in [y^{m+s}ax^nb]$.
  We note that, by the north-south dynamics of $y$, $w(y^{m-s}a)$ is to the left of $w(y^m a)$, and
  $w(y^{-s} b^{-1})$ to the left of $w(b^{-1})$.
  \begin{claim}
  \begin{itemize}
  \item $(y^{m-s} a x^n b)^+ \in w(y^{m-s} a)$
  \item $(y^{m-s} a x^n b)^- \in w(b^{-1})$
  \end{itemize}
  \end{claim}
  \begin{proof}
  It follows by Lemma~\ref{lem:cornering}, taking as the elements $a,b$ in that Lemma, the elements $y^{m-s} a$ and $b$. Note that the assumptions of the Lemma are satisfied: indeed, $b^{-1} \cdot \wt{x}$ and $a \cdot \wt{x}$ are parallel, by the assumptions of Lemma~\ref{lem:initial} and north-south dynamics of $y$ guarantees then that $b^{-1} \cdot \wt{x}$ and $y^{m-s} a \cdot \wt{x}$ are parallel.
  \end{proof}
  We also need the following claim, which follows similarly by Lemma~\ref{lem:cornering}.
  \begin{claim}
  \begin{itemize}
  \item $(y^{m} a x^n b y^{s})^+ \in w(y^{m} a)$
  \item $(y^{m} a x^n b y^{s})^- \in w(y^{-s} b^{-1})$
  \end{itemize}
  \end{claim}
  
Combined, these mean that
  $y^m a x^{n} b y^s$ crosses $y^{m-s}a x^nb$ to the right.  We have, thus,
  \[
    [a x^{n} b y^{m+s}][y^{m-s}ax^nby^{m-s}]=[y^m ax^{n}b y^s][y^{m-s}ax^nb] \reducesto [(y^max^nb)^2]= [(ax^nby^m)^2]
  \]
  proving Equation~\eqref{eq:crossuniform_n}.
  Equation~\eqref{eq:crossuniform_n} follows similarly: indeed, $\wt{y}$ being parallel to $a \cdot \wt{x}$ and $a \cdot \wt{x}$ being parallel to $ab \cdot \wt{y}$ is equivalent, by acting by $a^{-1}$, to $a^{-1} \cdot \wt{y}$ being parallel to $\wt{x}$ and $\wt{x}$ being parallel to $b \cdot \wt{y}$.
  Thus, we can consider a normalized bands model centered on $\wt{x}$ this time to prove Equation~\eqref{eq:crossuniform_n}.
\end{proof}

%%% Local Variables:
%%% mode: latex
%%% TeX-master: "Intersections"
%%% End:
